% Editable starting point. This file is never regenerated by catalogue tools.
\documentclass[aspectratio=169]{beamer}
\usetheme{bitpolito-studio}
\usepackage{bitpolitostudiotemplates}

\title{Your presentation}
\author{Your name}
\institute{Your institute}
\date{}

\begin{document}

% Standard pages use the document metadata and section names.
% Default: classic blue. Alternatives: [variant=bitgeneration] or [variant=bitgeneration-photo].
\studiotitleframe
\studiotocframe

\section{First ideas}
\studioframe{One clear idea}{%
  Use this default layout for simple content.
  \begin{itemize}
    \item Keep your message concise.
    \item Combine text, illustrations and diagrams.
  \end{itemize}
}

\section{Examples}
\studioframe{Three useful concepts}{%
  \begin{studiocards}[columns=3]
    \studiocard{\alice[height=1cm]}{Learn}
    \studiocard{\cow[height=1cm]}{Build}
    \studiocard{\bob[height=1cm]}{Share}
  \end{studiocards}
}

% For a measured layout, use a catalogue command instead of a studioframe.
% Place it outside any surrounding frame or studioframe. For example, when
% the three-card variant is included in your project:
% \studiotemplate[variant=three,title={Your concepts},
%   caption-one={Learn},caption-two={Build},caption-three={Share},
%   icon-one={alice},icon-two={cow},icon-three={bob}]{icon-cards}

\closingframe
\end{document}
